%% ============================================================================
%% §9 Conclusions  --  joint
%% ============================================================================
\section{Conclusions}
\label{sec:conclusion}

This project delivered a MobileFaceNet convolution frontend accelerator along two complementary paths: an ASIC in FreePDK45 and an FPGA on Zynq, both verified against the same C-model golden reference. Table~\ref{tab:final_summary} consolidates the headline metrics against the original constraints from Section~\ref{sec:constraints}.

\begin{table}[H]
\centering
\caption{Final Metrics vs Design Constraints}
\label{tab:final_summary}
\begin{tabularx}{\textwidth}{l c c c c}
\toprule
\headrow \textbf{Metric} & \textbf{Constraint} & \textbf{ASIC} & \textbf{FPGA} & \textbf{Met?} \\
\midrule
Clock frequency       & $\geq$100\,MHz (200 stretch) & 200\,MHz       & & \checkmark / \checkmark \\
Throughput            & $\geq$10\,FPS (target)       & 14.5\,FPS      & & \checkmark \\
Power                 & $<$100\,mW                   & 96.3\,mW       & & \checkmark \\
Core area             & $<$1\,mm$^2$                 & 0.518\,mm$^2$  & --      & \checkmark \\
DRC / antenna         & 0 violations                 & 0 / 0          & --      & \checkmark \\
50-layer MATCH        & required                     & 50/50 (RTL+syn+pnr) & & \checkmark \\
Timing signoff        & WNS $>$ 0                    & +0.089\,ns     & & \checkmark \\
\bottomrule
\end{tabularx}
\end{table}

\paragraph{Did the design meet spec?} Yes. All hard constraints were met; both soft constraints (power, area) cleared with margin. The ASIC path delivered 2$\times$ the stretch frequency target and stayed under the 100\,mW envelope. The functional correctness goal (50/50 layer MATCH against the C-model) was satisfied at every implementation stage, including post-route with SDF backannotation.

\paragraph{What was not delivered.}
\begin{itemize}[nosep]
\item No tape-out: the ASIC is design-kit-only.
\item No formal LEC: we used full GL simulation instead.
\item Single typical-PVT corner; no SI / cross-talk analysis.
\end{itemize}

\paragraph{Future work.}

\begin{table}[H]
\centering
\caption{Suggested Follow-on Work}
\label{tab:future}
\begin{tabularx}{\textwidth}{l X}
\toprule
\headrow \textbf{Priority} & \textbf{Task} \\
\midrule
High   & Multi-corner timing closure (slow / fast / typ); add signal-integrity-aware STA when tools allow. \\
High   & Formal LEC pass (Conformal / Formality) against RTL: seconds of run time replace days of GL simulation. \\
Medium & Integrate a small DMA / AXI shim and tape out a test chip on an MPW shuttle (FreePDK45 educational; real foundry would substitute its own PDK). \\
Medium & Replace combinational weight ROM with an external high-bandwidth weight SRAM + AXI-Lite write port for runtime model updates. \\
Low    & Mixed-precision exploration: INT8 activations + INT16 weights could roughly halve switching power while preserving the cosine-similarity gap. \\
% TODO friend: any FPGA-specific follow-on?
\bottomrule
\end{tabularx}
\end{table}

\paragraph{Lessons learned.} Verification took longer than RTL design. Two specific items dominated the schedule: (i) the gate-level testbench race that required \texttt{bind}-based probes to localise (Section~\ref{sec:opt:asic}), and (ii) the OpenRAM LEF patching needed to clean up DRC. Both would have been easier with deterministic instrumentation in place from day one.
